% Part: applied-modal-logics
% Chapter: epistemic-logic
% Section: properties-accessibility

\documentclass[../../../include/open-logic-section]{subfiles}

\begin{document}

\olfileid{aml}{el}{acc}

\olsection{অভিগম্যতা-সম্পর্ক ও জ্ঞানতাত্ত্বিক নীতিসমূহ}

স্বাভাবিক মোডাল যুক্তিবিদ্যায় ফ্রেম-সঙ্গতি সম্পর্কে আমরা যা জেনেছি, তার ভিত্তিতে জ্ঞানতাত্ত্বিক ব্যাখ্যায় পরিচায়ক !!{formula}s কেমন হয় তা দেখতে পারি। আগেই বলেছি, জ্ঞানতাত্ত্বিক যুক্তিবিদ্যার ব্যাখ্যায় সাধারণত S5 নেওয়া হয়। এখন বিভিন্ন জ্ঞানতাত্ত্বিক নীতি কীভাবে প্রকাশিত হয় এবং ফ্রেমের বিভিন্ন শর্তের সঙ্গে তাদের কী সঙ্গতি, তা দেখা যাক।

স্মরণ করুন, স্বাভাবিক মোডাল যুক্তিবিদ্যায় বিভিন্ন মোডাল !!{formula}s অভিগম্যতা-সম্পর্কের বিভিন্ন ধর্ম চিহ্নিত করেছিল। নিচের সারণিতে কয়েকটি সূত্র বেছে নেওয়া হয়েছে, যেগুলি নির্দিষ্ট জ্ঞানতাত্ত্বিক নীতির সঙ্গে মেলে।

\begin{table}[t]
    \begin{tabular}{| p{.48\textwidth} || p{.48\textwidth} |}
      \hline
      {\emph{$R$ যদি হয় \dots}} & {\emph{তবে~$\mModel{M}$-এ \dots সত্য:}} \\
      \hline \hline

      & $\Knows (p \lif q) \lif (\Knows p \lif \Knows q)$
      \hfill \newline{(নিহিতার্থে আবদ্ধতা)} \\
      \hline
      \emph{প্রতিবিম্বী}: $\forall w Rww$
      & $\Knows p \lif p$ \hfill \newline{(সত্য-অনুগামিতা)} \\
      \hline
      \emph{পরিযায়ী}: \newline
      $\forall u \forall v \forall w ((Ruv \land Rvw) \lif Ruw)$ &
      $\Knows p \lif \Knows \Knows p$ \hfill
           \newline{(ইতিবাচক অন্তর্দর্শন)} \\
      \hline
      \emph{ইউক্লিডীয়}: \newline
      $\forall w \forall u \forall v ((Rwu \land Rwv) \lif Ruv)$ &
      $\lnot \Knows p \lif \Knows \neg \Knows p$ \hfill
        \newline{(নেতিবাচক অন্তর্দর্শন)} \\
      \hline
    \end{tabular}
    \caption{চারটি জ্ঞানতাত্ত্বিক নীতি।}
    \ollabel{tab:four}
  \end{table}

$T$ স্বতঃসিদ্ধের সঙ্গে সংশ্লিষ্ট সত্য-অনুগামিতাকে প্রায়ই এই নীতিগুলির মধ্যে সবচেয়ে কম বিতর্কিত ধরা হয়। এর দাবি, কোনো !!a{formula} জানা থাকলে তা সত্যও হতে হবে। নিহিতার্থে আবদ্ধতা এবং ইতিবাচক ও নেতিবাচক অন্তর্দর্শন অনেক বেশি বিতর্কিত।

$K$ স্বতঃসিদ্ধের সঙ্গে সংশ্লিষ্ট নিহিতার্থে আবদ্ধতা বলে, কোনো কর্তার জ্ঞান নিহিতার্থের অধীনে আবদ্ধ। কোনো কোনো ক্ষেত্রে এটি বিশ্বাসযোগ্য মনে হতে পারে। যেমন, আমি জানি যে আমি যদি ভিক্টোরিয়ায় থাকি, তবে আমি ভ্যাঙ্কুভার দ্বীপে আছি। অস্বাভাবিক সংশয়বাদী পরিস্থিতি বাদ দিলে, আমি এটাও জানি যে আমি ভিক্টোরিয়ায় আছি; তাই আমার ভ্যাঙ্কুভার দ্বীপে থাকার কথাও জানা উচিত। এই ক্ষেত্রে আমার জ্ঞানের যৌক্তিক আবদ্ধতা বেশ স্বাভাবিক মনে হতে পারে। কিন্তু আমরা সব সময় আমাদের জানা কথার সব ফল নিয়ে চিন্তা করি না; ফলে অন্য ক্ষেত্রে এই নীতি কম স্বাভাবিক ফল দিতে পারে।

ইতিবাচক অন্তর্দর্শনকে কখনো KK-নীতি বলা হয়। এটি কখনো এই দাবিতে প্রকাশ করা হয়: আমি কিছু জানলে, যে আমি তা জানি সেটিও জানি। এটি 4 স্বতঃসিদ্ধের জ্ঞানতাত্ত্বিক প্রতিরূপ। অনুরূপভাবে নেতিবাচক অন্তর্দর্শন বলে: আমি কিছু \emph{না} জানলে, যে আমি তা জানি না সেটিও জানি; এটি 5 স্বতঃসিদ্ধের প্রতিরূপ। উভয় নীতির বিরুদ্ধেই সাধারণ জীবনের উদাহরণ মনে হয় পাওয়া যায়—যেমন, আমি বাস্তবে কোনো কিছু জানি, অথচ তা জানি কি না সে বিষয়ে নিশ্চিত নই।

\end{document}
